\documentclass[11pt]{article}
\usepackage[margin=0.85in]{geometry}
\usepackage{amsmath,amssymb,amsthm,mathtools}
\usepackage{graphicx}
\usepackage[colorlinks=true,linkcolor=blue,urlcolor=blue,citecolor=blue]{hyperref}
\usepackage{microtype}

\newtheorem{theorem}{Theorem}
\newtheorem{corollary}[theorem]{Corollary}
\newcommand{\Hh}{\mathcal H}
\newcommand{\Gg}{\mathcal G}
\newcommand{\Lat}{\operatorname{Lat}}

\title{Invariant Graph Subspaces Missing from the Claimed\\
Lattices of Powers of Unicellular Weighted Shifts}
\author{Candidate counterexample packet for arXiv:2205.00252}
\date{11 August 2026}

\begin{document}
\maketitle

\begin{center}
\fbox{\parbox{0.92\linewidth}{\textbf{Status: candidate counterexample likely valid.}
The construction gives a negative answer to the coordinate-form persistence
question in Remark 4.2.  More strongly, it supplies continuum many
counterexamples to Theorems 3.9--3.11 under the papers' own weight
hypotheses.}}
\end{center}

\section{Source and exact question}

Lata, Pokhriyal, and Singh study powers of unicellular weighted shifts in
\cite{LPS}.  In Remark 4.2 (arXiv PDF page 20; journal page 257), they ask
whether the lattice descriptions in Theorems 3.9--3.11 remain the same for
other sufficient classes of unicellular weights.

\begin{center}
\includegraphics[width=\linewidth]{figures/open_problem_crop.png}
\end{center}

In particular, Theorem 3.10 claims that, under the monotone square-summable
weight condition (3.4), every proper infinite-dimensional
$T^{*2}$-invariant subspace has one of a countable collection of coordinate
forms.  Theorem 3.11 gives an analogous coordinate description for
$T^{*3}$.

\section{Definitions}

Let $\Hh=\ell^2(\mathbb N_0)$ with orthonormal basis $(e_n)_{n\geq0}$.  For a
bounded sequence of positive weights $(w_n)$, let $B=T^*$ be the backward
weighted shift
\[
 Be_0=0,\qquad Be_n=w_{n-1}e_{n-1}\quad(n\geq1).
\]
For an integer $m\geq2$ and $0\leq r<m$, define the residue subspace
\[
 \Hh_r^{(m)}=\overline{\operatorname{span}}\{e_{mk+r}:k\geq0\}.
\]
Then $B^m\Hh_r^{(m)}\subseteq\Hh_r^{(m)}$, while
$B\Hh_r^{(m)}\subseteq\Hh_{r-1}^{(m)}$ when $1\leq r<m$.

\section{Idea of the proof}

The power $B^m$ preserves each residue class, but it also commutes with $B$.
Consequently the graph of the restriction
$B:\Hh_r^{(m)}\to\Hh_{r-1}^{(m)}$ is $B^m$-invariant.  The two residue
spaces are orthogonal, so this graph is automatically closed and cannot be a
coordinate subspace.  Varying its scalar slope produces continuum many such
graphs.  This elementary family is absent from the source paper's lists.

\section{Universal invariant-graph construction}

\begin{theorem}\label{thm:graphs}
Let $B$ be any backward weighted shift with positive bounded weights.  Fix
$m\geq2$, $1\leq r<m$, and $\lambda\in\mathbb C$.  Then
\[
 \Gg_{\lambda,m,r}:=(I+\lambda B)\Hh_r^{(m)}
\]
is a closed infinite-dimensional $B^m$-invariant subspace.  If
$\lambda\neq0$, it is proper and
\[
 \Gg_{\lambda,m,r}\cap\Hh_r^{(m)}
 =\Gg_{\lambda,m,r}\cap\Hh_{r-1}^{(m)}=\{0\}.
\]
For fixed $m,r$, distinct values of $\lambda$ give distinct subspaces.
\end{theorem}

\begin{proof}
For $x\in\Hh_r^{(m)}$, the vectors $x$ and $Bx$ belong to orthogonal residue
spaces.  Hence
\[
 \|(I+\lambda B)x\|^2=\|x\|^2+|\lambda|^2\|Bx\|^2\geq\|x\|^2.
\]
Thus $I+\lambda B$ is bounded below on $\Hh_r^{(m)}$, so its range there is
closed and infinite-dimensional.  Since $B$ commutes with $B^m$ and
$B^m\Hh_r^{(m)}\subseteq\Hh_r^{(m)}$,
\[
 B^m(I+\lambda B)x=(I+\lambda B)B^m x\in\Gg_{\lambda,m,r}.
\]

Assume $\lambda\neq0$.  The restriction of $B$ to $\Hh_r^{(m)}$ is
injective because $r\geq1$ and all weights are positive.  If
$x+\lambda Bx$ lies in $\Hh_r^{(m)}$, its orthogonal
$\Hh_{r-1}^{(m)}$ component gives $Bx=0$, hence $x=0$.  If it lies in
$\Hh_{r-1}^{(m)}$, its $\Hh_r^{(m)}$ component gives $x=0$.  The two stated
intersections are therefore trivial, and in particular the graph is proper.

Finally, if $\Gg_{\lambda,m,r}=\Gg_{\mu,m,r}$, comparison of the
$\Hh_r^{(m)}$ components shows that a vector $x$ must have the same graph
representative on both sides.  Thus $(\lambda-\mu)Bx=0$ for every $x$, and
injectivity gives $\lambda=\mu$.
\end{proof}

\section{Counterexamples to the source statements}

Take the Donoghue weights
\[
 w_n=2^{-n}.
\]
They are positive, monotonically decreasing, and square summable, so they
satisfy Condition (3.4) of \cite{LPS}.  They also satisfy Condition (3.1):
for a product of length $L\geq1$, the series occurring there is
$\sum_{k\geq0}2^{-2kL}\leq(1-2^{-2})^{-1}$.  Set
\[
 \Hh_{\mathrm{odd}}=\overline{\operatorname{span}}\{e_{2k+1}:k\geq0\},
 \qquad \Gg=(I+B)\Hh_{\mathrm{odd}}.
\]
By Theorem \ref{thm:graphs}, $\Gg$ is closed, proper,
infinite-dimensional, and $B^2$-invariant.  It meets both the even and odd
coordinate subspaces only at zero.  It is therefore neither coordinate
residue space in Theorems 3.9(i) and 3.10(i).  Every space in parts (ii) of
those theorems contains
$e_0$, whereas $\Gg$ does not.  Thus $\Gg$ is absent from the claimed
exhaustive lists.  Varying the nonzero scalar $\lambda$ gives continuum many
missing spaces, while each displayed list is countable.

The same argument with $m=3$ and $r=1$ contradicts Theorem 3.11: the graph
has zero intersection with each residue summand occurring in it, whereas
every nonzero coordinate direct sum in that theorem contains a pure residue
vector.  The forms in Theorem 3.11(iii) additionally contain $e_0$.

Consequently the premise behind Remark 4.2 is already false under one of the
paper's two original weight classes.  In particular, its coordinate-form
persistence question has a negative answer.

\section{The lattices genuinely depend on the weights}

The preceding objection is not merely that the common lists are incomplete.
As literal collections of subspaces, the lattices can differ even among
geometric weights satisfying Condition (3.4).

For $0<q<1$, let
\[
 B_qe_n=q^{n-1}e_{n-1}\quad(n\geq1),\qquad
 \Gg_q=(I+B_q)\Hh_{\mathrm{odd}}.
\]
Then $\Gg_q\in\Lat(B_q^2)$.  Let $p\in(0,1)$ and consider
$g=e_3+q^2e_2\in\Gg_q$.  Direct calculation gives
\[
 B_p^2g=p^3e_1+q^2p\,e_0.
\]
If this vector belonged to $\Gg_q$, its odd component would force its graph
preimage to be $p^3e_1$, whose even component is
$B_q(p^3e_1)=p^3e_0$.  Thus membership would require $q^2p=p^3$, or
$q=p$.  Taking $q=1/2$ and $p=1/3$ proves
\[
 \Gg_{1/2}\in\Lat(B_{1/2}^2)
 \quad\hbox{but}\quad
 \Gg_{1/2}\notin\Lat(B_{1/3}^2).
\]
Both weight sequences are within Condition (3.4).

\section{Verification and limitations}

The script \texttt{code/verify\_counterexample.py} uses exact rational
matrices.  It checks the graph construction for all residues and powers
$2\leq m\leq6$, distinct slopes, the square counterexample, and the
$q=1/2$ versus $p=1/3$ failure.  Its output is:
\begin{verbatim}
all exact finite-section checks passed
checked graph invariance for powers 2 through 6
verified the q=1/2 graph is not invariant for the p=1/3 square
\end{verbatim}
The computation is a guard against indexing errors; the proof is the
infinite-dimensional argument above.

A bounded novelty search on 11 August 2026 used the exact title, arXiv id,
authors, theorem number, ``lattice structure'', and combinations of
``weighted shift'', ``power'', and ``invariant graph subspace''.  It found
the arXiv version and the 2023 journal publication, but no correction,
erratum, or later source explicitly recording this counterexample.  This is
not a certification of novelty.

\paragraph{Human-review recommendation.}
Check the graph invariance identity and compare the two trivial-intersection
properties directly with Theorems 3.10 and 3.11.  Because the conclusion
contradicts published characterization statements, external use should wait
for independent expert review.

\begin{thebibliography}{9}
\bibitem{LPS}
S. Lata, S. Pokhriyal, and D. Singh,
\emph{Invariant subspaces of powers of some unicellular operators},
New York J. Math. \textbf{29} (2023), 231--260;
\href{https://arxiv.org/abs/2205.00252}{arXiv:2205.00252}.
\end{thebibliography}

\end{document}
